\section{The Appearance Ceiling}
\label{sec:ceiling}

\begin{table*}[t]
\centering
%    10 or 12-token rows ... The prefixed comment must stay stable.
% (see verify_claims.py parse_results_table strictness notes)
\caption{\textbf{Main result}, test split, $n_Q = 280$ query instances over 11
walls. All methods see identical inputs and are scored by identical code.
$F_1^{*}$ is oracle-thresholded on test; $\Delta F_1^{*}$ is its excess over the
$0.305$ best-constant-predictor floor ($2p/(1+p)$, not coin-flip chance;
Sec.~\ref{sec:protocol}); $F_1^{@v}$ is
thresholded on the held-out validation walls and is the deployable number.
\emph{scored} is the fraction of valid query--gallery pairs receiving a finite
score; unscored pairs are charged as predicted-negative and the query is still
counted (Sec.~\ref{sec:coverage}). Cost is wall-clock seconds per query. The
\textsc{Coverage-only} reference row answers yes to every registered pair
without examining a crack and is the level every other row must be read against;
chance levels for every column are given in Sec.~\ref{sec:protocol}. Best
crop-scope value per column in \textbf{bold}.}
\label{tab:main}
\small
\begin{tabular}{llcccccccr}
\toprule
Method & Family & R@1 & mAP & DIR@FAR$=$0.1 & $F_1^{*}$ & $\Delta F_1^{*}$ & $F_1^{@v}$ & scored & s/query \\
\midrule
\rowcolor{black!8}
\multicolumn{10}{l}{\emph{Reference level}} \\
\rowcolor{black!8}
\textsc{Coverage-only} & geometry      & 0.515 & 0.533 & 0.000 & 0.544 & 0.239 & 0.544 & 0.48 & --- \\
\midrule
\multicolumn{10}{l}{\emph{Crop scope}} \\
ORB~\cite{rublee2011orb}            & keypoint   & 0.352 & 0.364 & 0.018 & 0.305 & 0.000 & 0.305 & 1.00 & \best{0.038} \\
SIFT~\cite{lowe2004sift}            & keypoint   & 0.432 & 0.417 & 0.036 & 0.313 & 0.008 & 0.305 & 1.00 & 0.086 \\
SuperGlue~\cite{sarlin2020superglue}& learned    & 0.562 & 0.455 & \best{0.429} & 0.353 & 0.048 & 0.305 & 1.00 & 5.826 \\
LoFTR~\cite{sun2021loftr}           & learned    & 0.549 & 0.440 & 0.336 & \best{0.359} & \best{0.054} & \best{0.351} & 1.00 & 8.639 \\
YOLOv8 backbone                     & embedding  & 0.561 & 0.453 & 0.157 & 0.357 & 0.052 & 0.308 & 1.00 & 0.217 \\
ViT-B/16~\cite{dosovitskiy2021vit}  & embedding  & \best{0.654} & 0.511 & 0.286 & 0.331 & 0.026 & 0.317 & 1.00 & 0.193 \\
DeiT-S~\cite{touvron2021deit}       & embedding  & 0.621 & 0.497 & 0.250 & 0.305 & 0.000 & 0.293 & 1.00 & 0.058 \\
CLIP ViT-B/32~\cite{radford2021clip}& embedding  & 0.604 & 0.471 & 0.143 & 0.307 & 0.002 & 0.306 & 1.00 & 0.050 \\
OSNet~\cite{zhou2019osnet}          & re-ID      & \best{0.654} & \best{0.520} & 0.168 & 0.309 & 0.004 & 0.309 & 1.00 & 0.053 \\
\midrule
\rowcolor{black!5}
\multicolumn{10}{l}{\emph{Full-image scope}} \\
\rowcolor{black!5}
\ours{} (ours) & geometry & \best{0.952} & \best{0.886} & \best{0.768} & \best{0.656} & 0.351 & 0.582 & 0.48 & 27.92 \\
\bottomrule
\end{tabular}
\end{table*}

\subsection{Setup}
Nine crop-scope matchers --- handcrafted keypoints
(SIFT~\cite{lowe2004sift}, ORB~\cite{rublee2011orb}); learned correspondence
(SuperGlue~\cite{sarlin2020superglue}, LoFTR~\cite{sun2021loftr}); general
embeddings (ViT-B/16~\cite{dosovitskiy2021vit}, DeiT-S~\cite{touvron2021deit},
CLIP ViT-B/32~\cite{radford2021clip}, YOLOv8); a re-identification architecture
(OSNet~\cite{zhou2019osnet}) --- run under an identical protocol, inputs and
scoring.

\subsection{Result: appearance buys almost nothing on the decision}
On the decision a record actually makes --- \emph{same crack or not}, at a
threshold --- no crop method clears chance by more than $0.054$ (excess over the
$0.305$ floor at the oracle threshold, Table~\ref{tab:main}: $0.000$--$0.054$,
two methods exactly at the floor, eight within $0.052$), across two decades of
methods and two orders of compute ($\approx\!1$\,ms to $8.6$\,s). At a deployable (validation-calibrated)
threshold eight of the nine span only $-0.012$ to $+0.012$ around chance;
\textbf{LoFTR is the exception}, $F_1^{@v}=0.351$ ($+0.046$), and the most
expensive crop method. Ranking does not transfer to deciding: the nine separate
over a $0.302$ range on Rank-1, yet OSNet ties for best Rank-1 ($0.654$) while
sitting $0.004$ above chance, and SuperGlue (fifth on Rank-1) has the best crop
open-set rate ($0.429$). Ordering a gallery well and keeping a record are not
the same capability. The collapse after validation calibration is consistent
with threshold instability under a small, wall-clustered sample: methods can
order candidates differently while their score distributions do not support a
stable same-crack/not-same-crack boundary on held-out walls. This is an
interpretation of the observed pattern, not a claim that calibration itself is
the sole cause.

\subsection{Is it the crack, or the wall behind it?}
\label{sec:controls}
Two controls separate identity cue from field of view: a \emph{crop$+$context}
condition (wider crop) and a \emph{crack-erased} condition (the crack inpainted
out of that widened view). Context helps a lot --- OSNet's Rank-1 rises
$0.654 \to 0.775$, mAP to $0.595$, open-set rate $0.168 \to 0.404$ (CLIP's
$0.143 \to 0.311$; $F_1^{*}$ $0.309 \to 0.354$) --- \textbf{but erasing the
crack costs almost none of it}: OSNet retains Rank-1 $0.711$, mAP $0.574$ and
open-set $0.393$; CLIP retains Rank-1 $0.600$, mAP $0.502$, open-set $0.286$.
Both erased views still outrank every tight-crop method: the retrieved cue is
the wall, i.e.\ location. Pairwise $F_1$ is unmoved by context
($0.346$--$0.361$). We read the aggregate as a property of the data --- methods
with no shared inductive bias land near the floor, and wherever a crop method
improves, the
improvement is not the crack. Nine methods are a sample, not proof; the result
establishes that the obvious moves buy $0.000$--$0.054$ above chance on the
decision --- reason to change the formulation, not the encoder.